\section{Related Work}

\paragraph{Institutional holdings, herding and skill.}
A standard way to see what professional investors trade is to look at their
quarterly holdings. Many studies ask whether institutions herd, i.e.\ buy and
sell the same stocks at the same time~\cite{lakonishok1992impact,
wermers1999herding, sias2004herding}. Others ask whether institutions show
evidence of skill~\cite{chen2000value}, typically by comparing performance
with that of characteristic-matched benchmark
portfolios~\cite{daniel1997measuring}. For hedge funds, 13F filings show that
they participated in the technology bubble and then got out before it
burst~\cite{brunnermeier2004bubble}. The 13F filings provide only weak
evidence that hedge funds pick better stocks than mutual
funds~\cite{griffin2009smart}. Confidentially treated positions, whose
disclosure is delayed, earn abnormal returns~\cite{agarwal2013hide,
aragon2013confidential}; thus, the best informed positions may also be those
least visible. Our aim is to address a different question: how much of the
disclosed trades could have been predicted prior to disclosure? Our study
complements work on demand-system models of institutional
holdings~\cite{koijen2019demand}. Disclosed holdings may also be copied: by
mimicking the disclosed portfolios of other funds, funds can earn returns
similar to, if not higher than, those earned by the original
funds~\cite{frank2004copycat, verbeek2013copycat}. This makes the question of
how predictable disclosed trades are a practically relevant one as well.

\paragraph{Predicting fund managers' trades.}
Cohen, Lu and Nguyen~\cite{cohen2026mimicking} show that they can predict
71\% of active mutual fund managers' buy, hold and sell decisions, training an
LSTM per fund on Morningstar holdings data from 1990 through 2023. This beats
predicting each fund's most common action, which yields a 52\% accuracy.
Furthermore, they show that less predictable managers and positions perform
better. Izadifar and Rahmati~\cite{izadifar2026holdings} forecast the
portfolio weights of 99 managers using 13F data and temporal graph models.
They find that a simple persistence forecast beats all but one of their graph
models.

\paragraph{Foundation models in finance.}
Chronos~\cite{ansari2024chronos}, Chronos-2~\cite{ansari2025chronos2} and
TimesFM~\cite{das2024timesfm} are all capable of forecasting time series that
they have never seen before (i.e.\ they can forecast zero-shot). Rahimikia,
Ni and Wang~\cite{rahimikia2025tsfm} show that off-the-shelf foundation models
forecast daily excess returns poorly, and that models trained only on
financial data perform better. They release one such model per year, so we
can use one trained solely on data before our test period.
Noguer i Alonso and Pereira Franklin~\cite{noguer2026tsfm} test several
pre-trained models, including Chronos-2, on five large U.S.\ stocks and find
that their gains over a random walk are small and rare. These studies forecast prices, while we
test foundation models on the size of an investor's position.

\paragraph{News, sentiment and lookahead bias.}
We score headlines from FNSPID~\cite{dong2024fnspid} with
FinBERT~\cite{araci2019finbert} and use the policy-uncertainty index
of~\cite{baker2016epu}. Forecasts from large language models can be inflated
when the model has memorized what happened; Gao, Jiang and
Yan~\cite{gao2025lookahead} propose a test for this lookahead bias.
